\section{Discussion}
\subsection{When alternative policies should help}
\paragraph{Length and teacher specialization.}
Student confidence measures the student's uncertainty. The relevant quantity for routing is the teacher's chance of fixing an error minus its chance of replacing a correct answer with an error, conditional on what the gate can see. Length can improve this estimate if reviews of different lengths show different teacher benefit \emph{at comparable student margins}. A teacher trained for a particular genre, longer context, or a specific error type could make that difference large enough to matter. The present long-text average supports checking for such a signal. It cannot distinguish length from differences in difficulty or composition. In a separate diagnostic, the least-confident fifth of Stanford had teacher net gains of 1.54, 15.14, and 13.27 percentage points at up to 16, 17--32, and 33--64 tokens respectively. The middle bin led the longest of these bins even within this broad confidence group. Across the same diagnostic cohorts, the least-confident fifth gained about 11--12 points overall, while the remaining four fifths gained at most 0.35 points. That concentration explains why the margin was a strong baseline. To improve on it, new features must distinguish fixes from harms among requests that are comparably uncertain, especially near the cutoff.

\paragraph{Average cost and tail latency.}
For an additive mean-cost objective, suppose \(G(x)=\mathbb{E}[g\mid x]\) is expected net correctness gain, \(C(x)\) is expected incremental teacher time, and \(\lambda\) is the cost of one millisecond in correctness units. A fallback improves the objective when \(G(x)-\lambda C(x)>0\). Confidence remains competitive when this net value is already ordered by margin. Cost-aware routing can add value when teacher time varies substantially at similar margins, and when measured costs reflect the deployment workload. The power-law score in Equation~\ref{eq:score} is a heuristic for exploring this tradeoff. Summing per-bin medians estimates added work. It cannot determine p95, which depends on where slow requests land in the full latency distribution. A tail-sensitive deployment should measure latency by length and fallback outcome under representative load and select the policy under an explicit p95 constraint. Cutoffs should be calibrated on the target domain, especially when most of its requests collapse into one cost bin.

\paragraph{Input shape and student optimization.}
Fixed shapes are promising when requests cluster near a supported capacity and the measured specialization gain exceeds padding and routing overhead. Stable warm sessions and an affordable memory footprint are prerequisites for that comparison. The tested buckets offer no such demonstrated route. INT8 showed that a faster graph can still change the ranking of uncertain cases. Replacing FP32 would require an INT8-specific gate that meets the accuracy target on unseen labels, followed by matched timing. Cost-weighted training could help when the expensive group contains enough learnable student errors to justify shifting training emphasis. A cascade-aware training objective could instead prioritize the cases the student will actually keep. Neither quantization-aware training nor joint training of the student and gate was evaluated here.

\subsection{Limitations and scope of claims}
\textbf{Data and selection.} Gate fitting and calibration used SST-2 training examples from the public teacher's fine-tuning split. The partitions were separate from one another and from the student's training data. The teacher's training exposure can still distort estimates of its example-level benefit. The original operating points were frozen before row-level GLUE evaluation. Later sweeps used GLUE as development data, and Stanford was reused across the project and used to choose the 0.16 policy. A one-example gain after a large grid search does not establish superiority. The fresh IMDb audit excluded previously scored reviews, but it is label- and length-balanced and capped at 160 tokens. It does not represent all IMDb traffic, and differences between datasets cannot be attributed to length alone. A stronger replication would keep gate fitting and calibration out of teacher training, predeclare the policy and comparisons, and use a separate test set sampled for the target workload.

\textbf{Systems and training.} The CPU results use one laptop, batch-one requests, warm sessions, and a few paired process blocks. Repeated requests within a process are not independent hardware replications. I did not measure throughput under load, energy, or GPU performance. Training used one seed and had no four-layer supervised control, leaving the separate effects of depth and distillation unresolved at that size. I did not time the recalibrated INT8 gate, cost-weighted checkpoint, or external grid replays end to end.

\textbf{Claim boundary.} The accuracy budget is an observed engineering criterion. The original paired interval includes a loss larger than one point, so the study cannot establish population noninferiority. It also does not establish a universal length boundary, superiority over other models or tasks, or a novel cascade principle. Its contribution is the measured interaction between padding, model capacity, routing accuracy, and latency in this setting.

\section{Conclusion}
Selective fallback let the four-layer student meet the project's observed accuracy, latency, and memory goals on SST-2. The tuned cost-aware gate found a narrower advantage on that workload, exchanging a small median improvement and fewer calls for slower tail responses, with weaker transfer to longer IMDb reviews. Confidence-only is the better-supported default for these models. A targeted gate is most plausible when the deployment workload contains a reproducible source of teacher benefit beyond student margin, and when the saved computation is verified against the latency metric that the deployment actually cares about.
